\section{The computation of the Fourier coefficients}\label{app:computation}

The tables of this paper are Fourier coefficients of the inverse of the product \eqref{eq:product}, computed in exact rational arithmetic. We record the steps so that the numbers can be reproduced without the program, which is in the archive \cite{Bha26}.

\paragraphx{The genera.} The twisted elliptic genera are expanded in $q = e^{2\pi i\tau}$ and $y = e^{2\pi i z}$ from the theta functions $\theta_i(\tau,z)$ and $\eta(\tau)$ as $q$-series with coefficients in $\ZZ[y,y^{-1}]$, and the coefficients $c^{(r,s)}_b(4n-j^2)$ are read off from $\sum_{n,j} c^{(r,s)}_b(4n-j^2)\,q^n y^j$. For K3, $F^{(0,0)} = 8A$ with $A = \sum_{i=2}^4\theta_i(\tau,z)^2/\theta_i(\tau,0)^2$; for the $\ZZ_2$ CHL model the genera are $F^{(0,0)} = 4A$, $F^{(0,1)} = \tfrac43 A - \tfrac23 B\,E_2(\tau)$, $F^{(1,0)} = \tfrac43 A + \tfrac13 B\,E_2(\tau/2)$, $F^{(1,1)} = \tfrac43 A + \tfrac13 B\,E_2((\tau+1)/2)$, with $B = \theta_1(\tau,z)^2/\eta(\tau)^6$ and $E_2(\tau) = 2E(2\tau)-E(\tau)$, where $E$ is the Eisenstein series of weight two normalised to $1 - 24\sum_{n\ge1}\sigma_1(n)q^n$; for the $\ZZ_2$ toroidal orbifold they are \eqref{eq:tor2genera}. The arguments $\tau/2$ and $(\tau+1)/2$ are handled by expanding in $q^{1/2}$ with the sign $(-1)^n$ on the odd powers, so every series is a Laurent series in $x = q^{1/8}$, which is the smallest fractional power that occurs.

\paragraphx{The product.} Each factor $(1 - e^{2\pi i(k'\sigma + l\rho + jv)})^{-c}$ of $1/\Phit$ is expanded as a binomial series in the monomial $e^{2\pi i(k'\sigma + l\rho + jv)}$ to the order needed, and the factors are multiplied as polynomials in the three variables $e^{2\pi i\rho}$, $e^{2\pi i\sigma/N}$, $e^{2\pi iv}$ truncated at the exponents $Q^2/2\le L$, $NP^2/2\le K$ wanted in the tables. The exponent $c = \sum_s e^{-2\pi i s l/N}c^{(r,s)}_b(4k'l-j^2)$ is an integer. For $N\le 2$ the root of unity is $\pm1$; for $N=3,5,7$ the sum over $s$ is evaluated in closed form from the genera: for $r\ne0$ it is $\frac{8}{N+1}A$ at $N\mid l$ plus $\frac{2}{N+1}$ times a coefficient of $B(\tau)E_N(\tau/N)$, and for $r=0$ it is a rational combination of the coefficients of $A$ and of $B\,E_N$, so that no cyclotomic arithmetic is needed. The one factor with $k'=l=0$, $(1 - e^{\mp 2\pi iv})^{-2}$ after the prefactor is taken into account, is the only one whose expansion depends on the chamber: in the chamber $\operatorname{Im} v<0$ it is expanded in powers of $e^{-2\pi iv}$, in the chamber $\operatorname{Im} v>0$ in powers of $e^{2\pi iv}$. The two expansions give the two columns of Table~\ref{tab:phi10} at $\pm Q\cdot P$, and their difference is the wall-crossing jump of Theorem~\ref{thm:wallcrossing}, which the program checks for every $Q^2,P^2\in\{2,4,6\}$ and $|Q\cdot P|\le 4$.

\paragraphx{The index.} The coefficient of $e^{2\pi i(\rho Q^2/2 + \sigma P^2/2 + v\,Q\cdot P)}$ in $1/\Phit$, multiplied by $(-1)^{Q\cdot P+1}$, is the index $d(Q,P)$ of \eqref{eq:dvv}. For the orbifolds the prefactor of the product is $e^{2\pi i(\rho + \sigma/N + v)}$ for the CHL models and $e^{2\pi i(\rho + v)}$ for the toroidal orbifolds, as in \cite{CD19}, and the integral over $\sigma$ runs over a period of length $N$ (Remark~\ref{rem:normalisation}).

\paragraphx{The tables.} Table~\ref{tab:phi10} is the heterotic index for $Q^2,P^2\in\{2,4,6\}$ and $-2\le Q\cdot P\le 4$; it reproduces Table~1 of \cite{Sen11positivity} entry by entry. Table~\ref{tab:tor2} is the index of the $\ZZ_2$ toroidal orbifold for $Q^2\le 12$, $P^2\in\{0,2,4\}$ and $Q\cdot P\in\{0,1,2\}$; it agrees with Tables~17 to~19 of \cite{CD19} in every entry with $Q^2\le 12$ but one, the entry $Q^2=4$, $P^2=2$, $Q\cdot P=0$, which \cite{CD19} prints as $-11\,520$ and which the product formula gives as $+11\,520$, the sign required by the rule $(-1)^{Q^2}$ of \cite{CD19} for $P^2=2$; the misprint lies inside the single-centred cone and would otherwise be a thirty-first negative entry. Table~\ref{tab:violations} lists the charges in the cone \eqref{eq:z2conditions} with $Q^2\le 12$, $P^2\le 8$, $Q\cdot P\le 4$ whose index is negative. The same scan over the cone of the heterotic string and of the $\ZZ_2$ CHL model, up to $Q^2\le 12$ (heterotic) or $NQ^2/2\le 6$ (CHL), $P^2\le 6$, $Q\cdot P\le 4$, finds no non-positive index, and the CHL values reproduce Table~2 of \cite{Sen11positivity}.

\begin{table*}[t]
\centering\small
\setlength{\tabcolsep}{5pt}
\begin{tabular}{@{}cc*{7}{r}@{}}
\toprule
$Q^2$ & $P^2$ & \multicolumn{7}{c}{$d(Q,P)$ for $Q\cdot P = -2,\,-1,\,0,\,1,\,2,\,3,\,4$}\\
\midrule
2 & 2 & -209\,304 & 130\,329 & 50\,064 & 25\,353 & 648 & 327 & 0 \\
4 & 2 & -2\,023\,536 & 1\,598\,376 & 1\,127\,472 & 561\,576 & 50\,064 & 8\,376 & -648 \\
6 & 2 & -15\,493\,728 & 16\,844\,421 & 16\,491\,600 & 8\,533\,821 & 1\,127\,472 & 130\,329 & -15\,600 \\
4 & 4 & -16\,620\,544 & 28\,698\,000 & 32\,861\,184 & 18\,458\,000 & 3\,859\,456 & 561\,576 & 12\,800 \\
6 & 4 & -53\,249\,700 & 474\,507\,528 & 632\,078\,672 & 392\,427\,528 & 110\,910\,300 & 18\,458\,000 & 1\,127\,472 \\
6 & 6 & 2\,857\,656\,828 & 11\,890\,608\,225 & 16\,193\,130\,552 & 11\,232\,685\,725 & 4\,173\,501\,828 & 920\,577\,636 & 110\,910\,300 \\
\bottomrule
\end{tabular}
\caption{The index $d(Q,P) = (-1)^{Q\cdot P+1}c(Q,P)$ of quarter-BPS dyons of the heterotic string on $T^6$, from $1/\Phi_{10}$, in the chamber $\operatorname{Im} v<0$ (the expansion in powers of $e^{-2\pi i v}$). The table agrees entry by entry with Table~1 of \cite{Sen11positivity}. Entries with $Q\cdot P<0$ are the other chamber's values at $|Q\cdot P|$ by the symmetry $v\mapsto -v$; their negativity for $Q\cdot P\le -1$ is the two-centred contribution of Theorem~\ref{thm:wallcrossing}.}
\label{tab:phi10}
\end{table*}

\begin{table*}[t]
\centering\footnotesize
\setlength{\tabcolsep}{4pt}
\begin{tabular}{@{}c*{9}{r}@{}}
\toprule
& \multicolumn{3}{c}{$Q\cdot P=0$} & \multicolumn{3}{c}{$Q\cdot P=1$} & \multicolumn{3}{c}{$Q\cdot P=2$}\\
\cmidrule(lr){2-4}\cmidrule(lr){5-7}\cmidrule(lr){8-10}
$Q^2$ & $P^2=0$ & $P^2=2$ & $P^2=4$ & $P^2=0$ & $P^2=2$ & $P^2=4$ & $P^2=0$ & $P^2=2$ & $P^2=4$ \\
\midrule
0 & 2 & 64 & 816 & 0 & 0 & -8 & 0 & 0 & 0 \\
1 & -12 & -224 & -1\,248 & 1 & 96 & 1\,968 & 0 & 0 & -12 \\
2 & 48 & 1\,152 & 18\,240 & -8 & -256 & 840 & 0 & 64 & 2\,592 \\
3 & -168 & -3\,392 & -10\,320 & 37 & 1\,376 & 34\,656 & -2 & -224 & 2\,432 \\
4 & 528 & 11\,520 & 200\,736 & -136 & -3\,840 & 16\,632 & 16 & 1\,152 & 43\,392 \\
5 & -1\,512 & -30\,336 & -55\,424 & 439 & 13\,152 & 343\,152 & -72 & -3\,392 & 33\,720 \\
6 & 4\,032 & 83\,968 & 1\,544\,832 & -1\,288 & -33\,536 & 171\,152 & 256 & 11\,520 & 414\,336 \\
7 & -10\,176 & -202\,560 & -179\,022 & 3\,503 & 92\,928 & 2\,476\,752 & -806 & -30\,336 & 302\,400 \\
8 & 24\,528 & 496\,512 & 9\,480\,000 & -8\,968 & -220\,672 & 1\,265\,256 & 2\,320 & 83\,968 & 2\,926\,080 \\
9 & -56\,796 & -1\,118\,496 & -155\,232 & 21\,854 & 540\,416 & 14\,545\,584 & -6\,200 & -202\,560 & 2\,049\,968 \\
10 & 127\,008 & 2\,521\,600 & 49\,523\,328 & -51\,080 & -1\,204\,992 & 7\,558\,560 & 15\,616 & 496\,512 & 16\,919\,712 \\
11 & -275\,544 & -5\,374\,656 & 2\,684\,560 & 115\,154 & 2\,711\,616 & 73\,540\,080 & -37\,508 & -1\,118\,496 & 11\,568\,000 \\
12 & 581\,952 & 11\,389\,440 & 228\,872\,064 & -251\,528 & -5\,741\,824 & 38\,736\,600 & 86\,544 & 2\,521\,600 & 84\,554\,880 \\
\bottomrule
\end{tabular}
\caption{The index $d(Q,P)$ of quarter-BPS dyons in the $\ZZ_2$ toroidal orbifold of type IIB on $T^6$, from the product formula for $1/\Phit_2$ with the twisted elliptic genera \eqref{eq:tor2genera}, in the chamber $\operatorname{Im} v<0$ and with the normalisation of Remark~\ref{rem:normalisation}. The entries agree with Tables~17, 18 and~19 of \cite{CD19} in their common range $Q^2\le 12$, except the entry $Q^2=4$, $P^2=2$, $Q\cdot P=0$, printed there as $-11\,520$: the product formula gives $+11\,520$, and the sign rule $(-1)^{Q^2}$ for $P^2=2$ stated in the same section of \cite{CD19} requires the positive value. The column $P^2=0$ lies outside the single-centred cone; its sign is $(-1)^{Q^2}$. The negative entries inside the cone are those of Table~\ref{tab:violations}.}
\label{tab:tor2}
\end{table*}

